%% kapa: exam=1 round=19 year=2008 period=1 form=A subject=accounting subject_ko=회계학 seq=6 number=86
%% source: Q-Net 감정평가사 기출문제 (cst00309) artlSeq=1001962 file=extracted/1008/1001962/2008년 제19회 감정평가사 1차시험 기출문제/1교시.hwp
%% conversion: kordoc 4.24.51 -> tools/gen1.py (auto); review=proofread-sonnet; flags=none
\begin{question}{86}
다음 중 유형자산에 관한 기업회계기준서의 설명으로 적절하지 않은 것은? (단, 중소기업 회계처리 특례조항의 적용은 배제한다.)
\begin{choices}
  \item 유형자산을 장기후불조건으로 구입하거나, 대금지급기간이 일반적인 신용기간보다 긴 경우 취득원가는 취득시점의 현금구입가격으로 한다.
  \item 유형자산의 취득, 건설, 개발에 따른 복구비용에 대한 충당부채는 유형자산을 취득하는 시점에서 해당 유형자산의 취득원가에 반영한다.
  \item 유형자산을 자기회사가 직접 건설하는 경우 유형자산의 취득원가는 자가건설에 투입된 직접재료원가, 직접노무원가, 제조간접원가, 자가건설로 인한 생산감소에서 비롯되는 기회비용을 포함한다.
  \item 장기연불조건으로 매입한 유형자산의 취득원가는 자산취득으로 인해 생긴 부채의 명목가액과 현재가치의 차이가 중요한(material) 경우에는 현재가치로 평가한다.
  \item 시운전이나 본격적인 가동준비를 위한 지출이라도 유형자산을 사용가능한 상태로 만드는 과정과 직접적인 관련이 없다면 취득원가에 포함하지 않는다.
\end{choices}
\end{question}
\answer{3}
